\chapter{Estado de la Cuestión}
\label{cap:estadoDeLaCuestion}

\section{Conceptos biomédicos: RMN preclínica y lesiones cerebrales}

\section{Visión por computador en imagen biomédica}

La visión por computador o visión artificial es un campo de la inteligencia artificial que se centra en desarrollar métodos para transformar las imágenes
del mundo real en información que pueda ser interpretada por un ordenador. La visión por computador aplicada a imagen biomédica consiste en usar técnicas
automáticas para analizar imágenes biomédicas. Estas imágenes pueden venir de resonancia magnética, tomografía computarizada, ecografía o radiografía.
Cada modalidad muestra la información de una forma distinta y eso afecta al tipo de análisis que se puede hacer.

En medicina, muchas veces no basta con decir si una imagen tiene o no una lesión. También interesa saber dónde está, cuánto ocupa y qué parte del órgano
afecta. Por eso se usan tareas como la detección y la segmentación. Aunque la detección es una tarea parecida a la clasificación, la segmentación va un
paso más allá, ya que separa píxel a píxel la zona que se quiere estudiar del resto de la imagen.

Esta máscara de segmentación permite medir, por ejemplo, el área de una lesión en un corte o su volumen total si se trabaja con varios
cortes de una resonancia. Esto es útil porque convierte una observación visual en un dato numérico, no elimina la revisión experta, pero sí puede hacer
el proceso más rápido y más fácil de comparar entre casos.

En los últimos años, las redes neuronales convolucionales han tenido un gran uso en imagen biomédica. Estas redes aprenden patrones directamente desde las
imágenes, sin tener que definir a mano todas las características y pueden detectar bordes, cambios de intensidad o formas más complejas. En revisiones
como la de \citep{litjens2017survey} se recoge su uso en tareas de clasificación, detección y segmentación médica.

El problema es que en imagen médica es la falta de dato. Anotar imágenes lleva tiempo y suele requerir conocimiento experto. En segmentación,
además, no se etiqueta solo la imagen completa, sino que se dibuja la región de interés, esto hace que los conjuntos de datos sean más pequeños que en
otros campos.

También hay variaciones entre imágenes. Cambia el equipo, cambia el protocolo y cambia la modalidad. Incluso dos personas pueden segmentar de forma algo
distinta una misma lesión si el borde no se ve claro. Por eso, en este tipo de trabajos no basta con entrenar un modelo y mostrar una predicción. Hay que
comprobar si el resultado tiene sentido y si las medidas que se obtienen son coherentes.

En este Trabajo de Fin de Grado se aplica visión por computador al análisis de resonancias magnéticas preclínicas. El objetivo es automatizar parte de la
segmentación del cerebro y de la lesión, y después calcular medidas como el volumen.

\section{Arquitecturas de segmentación: U-Net}

La arquitectura U-Net, propuesta originalmente por Ronneberger et al. en 2015, fue diseñada para abordar las particularidades del análisis de imágenes biomédicas, un dominio con necesidades distintas a las de la visión por computador aplicada a imágenes naturales \citep{ronneberger2015unet}. En segmentación médica no basta con clasificar una imagen completa, sino que es necesario asignar una etiqueta a cada píxel o vóxel de la imagen de entrada, con el objetivo de delimitar con precisión estructuras anatómicas o regiones patológicas.

Además, los conjuntos de datos biomédicos anotados suelen ser reducidos, principalmente por restricciones de privacidad, disponibilidad limitada de datos y el elevado coste temporal y económico de la anotación experta. Por este motivo, los modelos destinados a este tipo de imágenes requieren arquitecturas eficientes en datos, capaces de aprender a partir de un número limitado de ejemplos.

U-Net presenta una estructura simétrica en forma de ``U'', compuesta por un camino contractivo o \textit{encoder} y un camino expansivo o \textit{decoder}. Esta organización responde a un objetivo funcional: procesar la información visual de forma jerárquica, combinando contexto semántico con precisión espacial.

El camino contractivo actúa como extractor jerárquico de características. En la formulación original de U-Net, este bloque aplica de forma repetida dos convoluciones \(3 \times 3\), seguidas de funciones de activación ReLU, y una operación de reducción espacial mediante \textit{max pooling} \(2 \times 2\). A medida que disminuye la resolución espacial de los mapas de características, el número de canales aumenta progresivamente. Por ejemplo, una imagen puede pasar de una resolución de \(128 \times 128\) a \(64 \times 64\), mientras que el número de filtros aumenta de 32 a 64, después a 128 y finalmente a 256 o 512. Este diseño permite ampliar el campo receptivo de la red y aprender representaciones cada vez más abstractas.

En las primeras capas, la red captura detalles locales de alta frecuencia, como bordes, texturas y gradientes de intensidad. Estas características son especialmente relevantes en resonancia magnética, ya que ayudan a definir los límites inmediatos de una región de interés. A medida que la red profundiza y la resolución espacial disminuye, el campo receptivo efectivo aumenta, permitiendo que las capas más profundas representen patrones de mayor nivel. En este punto, el \textit{bottleneck} concentra la información semántica de la imagen, es decir, el ``qué'' aparece en ella, aunque a costa de perder parte de la precisión espacial, es decir, el ``dónde''.

El camino expansivo trata de revertir este proceso para reconstruir el mapa de segmentación final. Para ello, aumenta progresivamente la resolución de los mapas de características hasta recuperar el tamaño espacial necesario para generar una máscara de salida. Sin embargo, un simple sobremuestreo de las características profundas puede producir contornos imprecisos, ya que parte de la información espacial fina se pierde durante las operaciones de reducción del \textit{encoder}. Esta limitación se aborda mediante uno de los elementos más característicos de U-Net: las conexiones de salto o \textit{skip connections}.

Las \textit{skip connections}, implementadas habitualmente mediante concatenación, permiten fusionar los mapas de características del \textit{encoder} con los mapas sobremuestreados del \textit{decoder} en niveles de resolución equivalentes. De esta forma, el \textit{decoder} recibe simultáneamente información semántica procedente de las capas profundas e información espacial de alta resolución procedente de las capas iniciales.

De manera simplificada, si \(x_l\) representa la salida del nivel \(l\) del \textit{encoder} e \(y_{l+1}\) representa la salida del nivel más profundo del \textit{decoder}, la entrada del nivel correspondiente del \textit{decoder} puede expresarse como:

\[
\text{Input}_{decoder,l} =
\operatorname{Concat}\left(\operatorname{Up}(y_{l+1}), x_l\right)
\]

donde \(\operatorname{Up}(\cdot)\) representa la operación de sobremuestreo y \(\operatorname{Concat}(\cdot)\) la concatenación de características. Esta combinación permite recuperar detalle espacial sin renunciar al contexto global aprendido por la red.

\section{Aprendizaje por transferencia (Transfer Learning)}

\section{Explicabilidad en Inteligencia Artificial (XAI)}